%% kapa: exam=1 round=35 year=2024 period=1 form=A subject=real_estate subject_ko=부동산학원론 seq=21 number=101
%% source: Q-Net 감정평가사 기출문제 (cst00309) artlSeq=5213229 file=raw/1008/5213229/1교시 문제지 원본.pdf
%% conversion: kordoc 4.24.51 -> tools/gen1.py (auto); review=proofread-sonnet; flags=none
\begin{question}{101}
X 노선 신역사가 들어선다는 정보가 있다. 만약 부동산 시장이 할당효율적이라면 투자자가 최대한 지불할 수 있는 정보비용의 현재가치는? (단, 제시된 가격은 개발정보의 실현 여부에 의해 발생하는 가격차이만을 반영하고, 주어진 조건에 한함)
\begin{framedtext}
○ X 노선 신역사 예정지 인근에 일단의 A 토지가 있다. \par
○ 1년 후 도심에 X 노선 신역사가 들어설 확률이 60 \%로 알려져 있다. \par
○ 1년 후 도심에 X 노선 신역사가 들어서면 A 토지의 가격은 5억 5,000만원, 신역사가 들어서지 않으면 2억 7,500만원으로 예상된다. \par
○ 투자자의 요구수익률(할인율)은 연 10 \%이다. \par
\end{framedtext}
\begin{choices}
  \item 5천만원
  \item 1억원
  \item 1억 5천만원
  \item 2억원
  \item 2억 5천만원
\end{choices}
\end{question}
\answer{2}
